\chapter{Difference Becomes Countable}

\section{Barely more than telling apart}

The ground floor held a single separation, $p \neq q$, and refused to call it a number. This chapter takes the one step that
turns separations into counts, and the step is smaller than it looks: to count is to gather distinguishables that share a
characteristic into a single heap and read off its size. Nothing new is assumed. The count falls out of the distinction the
moment there is more than one of it. Where the last chapter had two points told apart, this one has a collection of points,
each told apart from the rest, and the collection's size is the first number the geometry ever writes.

The physics where this is exact is the spectrum. Lay an operator $\hat{A}$ on the space and ask for its eigenvalues; on a
compact domain, or for a bound state, the answer is not a continuum but a discrete ladder, $\hat{A}\,\psi_n = a_n\,\psi_n$,
with $n = 0, 1, 2, \dots$ indexing distinguishable states. Planck reached the first such ladder when he found that the
oscillators of a black body could not carry arbitrary energies but only the discrete rungs $E_n = n\,h\nu$~\cite{planck1901,
planck1900}; the count $n$ was forced on him by the data, not chosen. What makes the spectrum countable is exactly what made
the two points distinguishable: each state is told apart from its neighbors by some characteristic --- a node, an energy, a
quantum number --- and the states that share none of the distinguishing marks are, for the count, one. The integer $n$ is
the size of the heap of states the resolution cannot split finer.

This is the first appearance of the second verb, and it is native here. To \emph{funge} is to gather by sameness --- to bag
into one covariant reading everything a chosen basis cannot tell apart. On this face a funge is the covariant act: the
component $V_\mu$ of a tensor is a single number standing for a whole equivalence class of the world's behavior along the
$\mu$-direction, everything the basis vector $e_\mu$ pairs the same way. To count the spectrum is to funge: to declare
$\psi_n$ and $\psi_{n}$ the same state whenever the instrument cannot resolve a difference, and to let the number of such
classes be the reading. Where tange laid a contravariant label on a point, funge gathers the world's covariant response into
a countable ladder. The count is a funge completed --- a heap gathered, sized, and named by its cardinality.

And the ladder inherits its integrity from the distinction beneath it. The reason $n$ is an \emph{integer} and not a smear is
that the states it indexes are genuinely told apart --- orthogonal, $\langle \psi_m \mid \psi_n \rangle = \delta_{mn}$, each
one distinguishable from every other by the inner product the space carries. Orthogonality is distinctness made quantitative:
two states are countably-two exactly when their overlap vanishes. So the count does not add a new principle to the ground
floor; it applies the ground floor's one principle, distinctness, to a collection, and reads the collection's size. The
Kronecker delta $\delta_{mn}$ is the funge written down --- one on the diagonal, where a state is bagged with itself, zero
off it, where the resolution keeps two states apart.

What the reader should carry forward is how little was spent. No continuum was invoked to reach the count; the ladder is
discrete because the distinctions are discrete, and a discrete tower of told-apart states is countable by its own
construction. The number is the size of a funge over tanged separations, and that is the entire content of ``countable'' at
this tier. The instrument has its first magnitude now --- but only an integer, and only the integer the resolution earns.
The reals are a further, harder object, and the next section builds them the honest way, as things approached and never
completed.

\section{A number is a finite condition}

The count of the last section was an integer, and integers are the easy numbers --- each one is exhibited whole, in a finite
act. The reals are not like that, and the geometry is honest about the difference. A real number, on this construction, is
never handed over complete; it is \emph{approached} through a tower of finite conditions, each one writable, none of them the
number itself. This section builds the real as what it is on an honest bench --- a limit of finite approximations, a Cauchy
tower --- and marks, once and for all, that the instrument holds the approximations and never the limit.

Take the definition at face value. A real $x$ is a Cauchy sequence of rationals $(x_1, x_2, x_3, \dots)$: for every
tolerance $\varepsilon > 0$ there is a stage $N$ past which $\lvert x_m - x_n \rvert < \varepsilon$ for all $m, n \geq N$.
Each $x_k$ is a finite, exhibited rational --- a thing you can write down and check. The real $x$ is not any of them; it is
the pattern they make as they crowd together, the address the tower converges on. Cauchy's condition is exactly a statement
about the \emph{approach} and says nothing that requires the destination to be handed over: it certifies that the terms grow
mutually close, that the tower is genuinely converging, without ever exhibiting a completed limit. This is the mathematics
the instrument runs on --- finite conditions, converging, the limit approached and never reached.

The physics reads this as the plain fact that every measured length is a finite approximation to a real one. A distance
$\mathrm{d}s = \sqrt{g_{\mu\nu}\,\mathrm{d}x^\mu\,\mathrm{d}x^\nu}$ is, in principle, a real; a bench returns it as a rational
to some number of places, with a next place always in principle available and never actually completed. The apparatus does
not fail to reach the real out of clumsiness. It is structurally a device that produces finite conditions --- a reading to
this tolerance, then a finer reading to that tolerance --- a Cauchy tower whose limit is the ideal value it climbs toward and
never occupies. The distinction between the tower and its limit is not a nicety. It is the exact seam along which, chapters
from now, the honest bracket will be cut: the instrument owns every rung and never the top.

Here the construction plants a fence it will not pull up. The continuum --- the completed set of all reals, each present
whole --- is precisely the object the finite instrument declines to assume. It builds reals as \emph{approached}, one finite
condition after another, and never as a completed totality handed down from above. This is the same fence the preface drew
around the axiom of choice: a completed continuum, with every real simultaneously in hand and selectable, is a tange promoted
past the finite, and the instrument refuses it. What it keeps is the tower: writable rungs, a certified approach, and a limit
named but never occupied. A great deal of the book's honesty is already latent in this refusal, because a machine that
helped itself to the completed continuum could pluck any real it liked from it and call the plucking a measurement.

In the two verbs, the tower is a funge that never closes. Each rung $x_k$ funges the value to a tolerance --- bags it within
$\varepsilon$ of its neighbors --- and the next rung funges it finer, but no rung completes the bag, because the real sits at
a limit no finite condition reaches. The instrument funges the value to arbitrary tightness and declines, at every stage, the
final tange that would pin it to a point. So a number, at this tier, is a finite condition and a promise of finer ones ---
never the completed real, always the rung and the approach. The count of the last section was a funge the instrument could
close, because integers are finite heaps; the real of this one is a funge it holds open forever, and that open funge is the
first shadow of the bracket the whole book is walking toward.

\section{The first loop}

The last two sections built a count and then a number-as-approach. This one asks the question that turns a number into a
measurement rather than a bookkeeping entry: does the number \emph{survive} being carried away and brought back? A reading
that changed every time you re-expressed it would be no reading at all. The test of a genuine count is that it comes back the
same after a round trip through the representations --- that it is invariant under the relabelings the instrument can
perform. This section runs that round trip, and finding the count unchanged at the end of it is the first loop, the
prototype of every measurement to come.

The mathematics of survival is invariance under a change of frame. Take the count and express it as a scalar --- a tensor
with no free index, $S$, whose transformation law under $x^\mu \mapsto x'^\mu$ is the trivial one, $S' = S$. A scalar is by
definition the quantity that every frame agrees on; it is what a covariant index and a contravariant index produce when they
are summed against each other, $S = V_\mu V^\mu = g_{\mu\nu} V^\mu V^\nu$, the frame-dependent labels cancelling to leave a
frame-free number. The count of distinguishable states is such a scalar. The dimension of an eigenspace, the number of
orthogonal states below a cutoff, the winding of a phase around a loop --- these are integers no change of coordinates can
alter, because they count distinctions the coordinates did not create and cannot destroy. Carry the count through any chart
and back; it returns unchanged, $S' = S$, and that return is the loop closing.

The reason the count survives is that it was built from distinctions, and distinctions are prior to the labels. A change of
frame reshuffles the contravariant tange --- our coordinates, our chosen basis --- but it cannot merge two orthogonal states
into one or split one into two, because orthogonality is a covariant fact about the world's inner product, not about our
labels. So the funge that produced the count is stable under exactly the operations that reshuffle the tange. This is the
first place the two indices are seen doing their opposite jobs in concert: the tange transforms, the funge counts, and their
contraction is the invariant that closes the loop. The scalar is the count that does not care which labels carried it.

There is a caution to state plainly, because the reals of the last section do not close this loop as cleanly as the integers
do. An integer count returns \emph{exactly}; a real, held only as a Cauchy tower, returns to within a tolerance --- the round
trip through representations preserves the approach but need not preserve the last digit, because the last digit was never
completed. So the first loop closes exactly on the count and only approximately on the real, and that gap --- exact on the
distinction, approximate on the magnitude --- is not a defect to be repaired. It is the same seam again, and the instrument
will read it, chapters on, as the width of a bracket. The loop that closes exactly is the structure; the loop that closes
only to a tolerance is where the number lives.

So the chapter ends with a number that is genuinely a reading: a count, funged from tanged distinctions, that survives its
own relabeling and returns invariant. That survival is the difference between a magnitude and a measurement --- a magnitude
is a value, a measurement is a value that comes back the same. The instrument has now taken a difference, gathered it into a
count, approached a real, and carried the whole through a loop that returned it unchanged. What it has not yet met is the
loop that \emph{fails} to close --- the round trip that comes back altered, and whose failure to return is not a flaw but the
signal the entire instrument is built to read. That failure is the residue, and it is the subject of the next chapter.
